Päivystyspoliklinikoiden ruuhkautuminen vaarantaa potilasturvallisuuden ja heikentää toiminnan tehokkuutta. Tarkka kuormituksen ennustaminen mahdollistaisi resurssien kohdentamisen ennakoivasti. Tässä työssä arvioidaan, pystyvätkö esikoulutetut aikasarjamallit (engl. foundational time series models) ennustamaan päivystyksen potilasmääriä ilman mallin erillistä kouluttamista historiallisen potilasdatan avulla (engl. zero-shot). Tämä kyvykkyys poistaisi perinteisten koneoppimismallien kouluttamiseen vaadittavat laskennalliset kustannukset ja manuaalisen työn.

Tässä työssä vertaillaan viittä avoimen lähdekoodin esikoulutettua mallia (TimesFM, Chronos, Sundial, Moirai ja TiRex) suorituskyvyltään edistyksellisimpiin koneoppimismalleihin käyttäen Tampereen yliopistollisen sairaalan päivystyksen kuormitusdataa. Vertailukelpoisuuden saavuttamiseksi esikoulutetut mallit arvioidaan käyttämällä samaa dataa kuin vertailumallien arvioimiseen käytettiin. Tutkimuksessa arvioitiin ennustetarkkuutta, historiatietojen pituuden ja ulkoisten muuttujien vaikutusta, sekä laskennallisia vaatimuksia.

Tulokset osoittavat, että yhden muuttujan perusmallit, erityisesti TimesFM, TiRex ja Sundial, suoriutuvat paremmin kuin paras aiempi vertailumalli, kun niille syötetään riittävästi historiallista aikasarjadataa. Aikasarjadatan pidentäminen paransi tarkkuutta johdonmukaisesti aina suurimpaan mahdolliseen 12 840 tunnin pituuteen asti. Ainoa monimuuttujadataa tukeva malli, Moirai, ei kyennyt hyödyntämään ulkoisia muuttujia. Laskennallisesti kaikkien mallien laskenta-ajat ja muistinkäyttö ovat kohtuullisella tasolla nykyaikaisille tietokoneille. Tulokset viittaavat siihen, että esikoulutetut aikasarjamallit tarjoavat varteenotettavan ja skaalautuvan vaihtoehdon perinteisille ennustusmenetelmille päivystyksen ruuhkautumisen hallinnassa.
